\section{AHSME 1959}
        \begin{enumerate}
        	\item (\textit{AHSME 1959}) Each edge of a cube is increased by $50$\%. The percent of increase of the surface area of the cube is:
$\textbf{(A)}\ 50 \qquad\textbf{(B)}\ 125\qquad\textbf{(C)}\ 150\qquad\textbf{(D)}\ 300\qquad\textbf{(E)}\ 750$


 
 \textbf{Answer: }B



\textbf{Solution}

Note that increasing the length of each edge by $50$\% with result in a cube that is similar to the original cube with scale factor $1.5$. Therefore, the surface area will increase by a factor of $1.5^2$, or $2.25$. Converting this back into a percent, the percent increase will be $125$\%. Therefore, the answer is $\boxed{\textbf{B}}$.

 


	\item (\textit{AHSME 1959}) Through a point $P$ inside the $\triangle ABC$ a line is drawn parallel to the base $AB$, dividing the triangle into two equal areas. 
If the altitude to $AB$ has a length of $1$, then the distance from $P$ to $AB$ is:
$\textbf{(A)}\ \frac12 \qquad\textbf{(B)}\ \frac14\qquad\textbf{(C)}\ 2-\sqrt2\qquad\textbf{(D)}\ \frac{2-\sqrt2}{2}\qquad\textbf{(E)}\ \frac{2+\sqrt2}{8}$


 
 \textbf{Answer: }D



\textbf{Solution}


                    \begin{figure}[H]
                        centering
                        \includegraphics[width=0.5\textwidth]{Solution/AHSME_I_1959_2_Solution_Figure_1.png}
                    \end{figure}
                    

 Because $[ABDE]=[\triangle EDC]$, we know that $[\triangle EDC]=\frac{1}{2}[\triangle ABC]$. Because the ratio of their areas is $\frac{1}{2}$, we know that the ratio of their side lengths (and thus their altitudes) is $\sqrt{\frac{1}{2}}=\frac{\sqrt{2}}{2}$. Thus, the altitude from $C$ to $\overline{DE}$ has length $1*\frac{\sqrt2}{2}=\frac{\sqrt{2}}{2}$. Thus, because $\overline{DE} \parallel \overline{AB}$, the distance from $P$ to $\overline{AB}$ is $1-\frac{\sqrt{2}}{2}=\boxed{\textbf{(D) }\frac{2-\sqrt{2}}{2}}$.

 


	\item (\textit{AHSME 1959}) If the diagonals of a quadrilateral are perpendicular to each other, the figure would always be included under the general classification:
$\textbf{(A)}\ \text{rhombus} \qquad\textbf{(B)}\ \text{rectangles} \qquad\textbf{(C)}\ \text{square} \qquad\textbf{(D)}\ \text{isosceles trapezoid}\qquad\textbf{(E)}\ \text{none of these}$


 
 \textbf{Answer: }E



\textbf{Solution}

Note that a \href{https://artofproblemsolving.com/wiki/index.php?title=Kite}{kite} has perpendicular diagonals, but is not a subcategory of any of the existing answer choices. Therefore, the answer is $\boxed{\textbf{E}}$.

 


	\item (\textit{AHSME 1959}) If $78$ is divided into three parts which are proportional to $1, \frac13, \frac16,$ the middle part is:
$\textbf{(A)}\ 9\frac13 \qquad\textbf{(B)}\ 13\qquad\textbf{(C)}\ 17\frac13 \qquad\textbf{(D)}\ 18\frac13\qquad\textbf{(E)}\ 26$


 
 \textbf{Answer: }C



\textbf{Solution}

Let the part proportional to $1$ equal $x$. Then, the parts are $x, \frac{1}{3}x$, and $\frac{1}{6}x$. The sum of these parts should be $78$, so $\frac{9}{6}x=1\frac{1}{2}x=78$. Solving for $x$, $x=52$. The middle part is $\frac{1}{3}x$, so the answer is $17\frac{1}{3}$, or $\boxed{\textbf{C}}$

 


	\item (\textit{AHSME 1959}) The value of $\left(256\right)^{.16}\left(256\right)^{.09}$ is:


 $\textbf{(A)}\ 4 \qquad \\ \textbf{(B)}\ 16\qquad \\ \textbf{(C)}\ 64\qquad \\ \textbf{(D)}\ 256.25\qquad \\ \textbf{(E)}\ -16$


 
 \textbf{Answer: }A



\textbf{Solution}

When we multiply numbers with exponents, we add the exponents together and leave the bases unchanged. We can apply this concept to computate $256^{0.16} \cdot 256^{0.09}$: 
 \[256^{0.16} \cdot 256^{0.09} = 256^{0.16+0.09}=256^{0.25}.\]Now we can convert the decimal exponent to a fraction: 
 \[256^{0.25} = 256^{\frac{1}{4}}.\]Now, let us convert the expression into radical form. Since $4$ is the denominator of the fractional exponent, it will be the index exponent: 
 \[256^{\frac{1}{4}}=\sqrt[4]{256}.\]Since $256 =16^2=(4^2)^2=4^4$, we can solve for the fourth root of $256$: 
 \[\sqrt[4]{256}=\sqrt[4]{4^4}=4.\]Therefore, $(256)^{.16} \cdot (256)^{.09}=\boxed{\textbf{(A) }\ 4}.$

 


	\item (\textit{AHSME 1959}) Given the true statement: If a quadrilateral is a square, then it is a rectangle. 
It follows that, of the converse and the inverse of this true statement is:


 $\textbf{(A)}\ \text{only the converse is true} \qquad \\ \textbf{(B)}\ \text{only the inverse is true }\qquad  \\ \textbf{(C)}\ \text{both are true} \qquad \\ \textbf{(D)}\ \text{neither is true} \qquad \\ \textbf{(E)}\ \text{the inverse is true, but the converse is sometimes true}$


 
 \textbf{Answer: }D



\textbf{Solution}

First, let us list the statement "If a quadrilateral is a square, then it is a rectangle" as a statement of the form "If $p$, then $q$". In this case, $p$ is "a quadrilateral is a square", and $q$ is "it is a rectangle".

 The converse is then: "If $q$, then $p$". Plugging in, we get "If a quadrilateral is a rectangle, then it is a square". However, this is obviously false, as a rectangle does not have to have four sides of the same measure.

 The inverse is then: "If not $p$, then not $q$". Plugging in, we get "If a quadrilateral is not a square, then it is not a rectangle". This is also false, since a quadrilateral can be a rectangle but not be a square. 

 Therefore, the answer is $\boxed{\textbf{D}}$.

 


 


	\item (\textit{AHSME 1959}) The sides of a right triangle are $a$, $a+d$, and $a+2d$, with $a$ and $d$ both positive. The ratio of $a$ to $d$ is:
$\textbf{(A)}\ 1:3 \qquad\textbf{(B)}\ 1:4 \qquad\textbf{(C)}\ 2:1\qquad\textbf{(D)}\ 3:1\qquad\textbf{(E)}\ 3:4$


 
 \textbf{Answer: }D



\textbf{Solution}

If we let $a=3$ and $d=1$, then we will get a $3$-$4$-$5$ triangle, which is a right triangle. So, the answer is $\boxed{\textbf{(D)} \ 3:1}$.

 


	\item (\textit{AHSME 1959}) The value of $x^2-6x+13$ can never be less than:


 $\textbf{(A)}\ 4 \qquad \textbf{(B)}\ 4.5 \qquad \textbf{(C)}\ 5\qquad \textbf{(D)}\ 7\qquad \textbf{(E)}\ 13$


 
 \textbf{Answer: }A



\textbf{Solution}

The $x$ value at which the minimum value of this \href{https://artofproblemsolving.com/wiki/index.php?title=Quadratic}{quadratic} (of the form $ax^2+bx+c$) occurs is $\frac{-b}{2a}=-\frac{-6}{2\cdot1}=3$. The minimum value of the quadratic is therefore at
 \[3^2-6\cdot3+13\]
 \[=9-18+13\]
 \[=4.\]So, the answer is $\boxed{\textbf{(A)} \ 4}$.

 


	\item (\textit{AHSME 1959}) A farmer divides his herd of $n$ cows among his four sons so that one son gets one-half the herd, 
a second son, one-fourth, a third son, one-fifth, and the fourth son, $7$ cows. Then $n$ is:


 $\textbf{(A)}\ 80 \qquad\textbf{(B)}\ 100\qquad\textbf{(C)}\ 140\qquad\textbf{(D)}\ 180\qquad\textbf{(E)}\ 240$


 
 \textbf{Answer: }C



\textbf{Solution}

The first three sons get $\frac{1}{2}+\frac{1}{4}+\frac{1}{5}=\frac{19}{20}$ of the herd, so that the fourth son should get $\frac{1}{20}$ of it. But the fourth son gets $7$ cows, so the size of the herd is $n=\frac{7}{\frac{1}{20}} = 140$. Then our answer is $\fbox{\textbf{(C) }140}$.

 


	\item (\textit{AHSME 1959}) In $\triangle ABC$ with $\overline{AB}=\overline{AC}=3.6$, a point $D$ is taken on $AB$ at a distance $1.2$ from $A$. 
Point $D$ is joined to $E$ in the prolongation of $AC$ so that $\triangle AED$ is equal in area to $ABC$. Then $\overline{AE}$ is:
$\textbf{(A)}\ 4.8 \qquad\textbf{(B)}\ 5.4\qquad\textbf{(C)}\ 7.2\qquad\textbf{(D)}\ 10.8\qquad\textbf{(E)}\ 12.6$


 
 \textbf{Answer: }D



\textbf{Solution}


                    \begin{figure}[H]
                        centering
                        \includegraphics[width=0.5\textwidth]{Solution/AHSME_I_1959_10_Solution_Figure_1.png}
                    \end{figure}
                    

 Note that $[\triangle ABC]=[\triangle ADE]$, so $\frac{1}{2}AB * AC *\sin\angle BAC = \frac{1}{2}AD * AE *\sin\angle DAE$. Since $\angle BAC = \angle DAE$, we have $AB*AC = AD*AE$, so that $3.6*3.6 = 1.2*AE$. Therefore, $AE = \frac{3.6^2}{1.2}=\boxed{\textbf{(D) }10.8}$.

 


	\item (\textit{AHSME 1959}) The logarithm of $.0625$ to the base $2$ is:
$\textbf{(A)}\ .025 \qquad\textbf{(B)}\ .25\qquad\textbf{(C)}\ 5\qquad\textbf{(D)}\ -4\qquad\textbf{(E)}\ -2$


 
 \textbf{Answer: }D



\textbf{Solution}


 \[\log_{2}{0.0625}=\log_{2}{\frac{1}{16}}=\boxed{-4}.\]

 


	\item (\textit{AHSME 1959}) By adding the same constant to $20,50,100$ a geometric progression results. The common ratio is:
$\textbf{(A)}\ \frac53 \qquad\textbf{(B)}\ \frac43\qquad\textbf{(C)}\ \frac32\qquad\textbf{(D)}\ \frac12\qquad\textbf{(E)}\ \frac{1}3$


 
 \textbf{Answer: }A



\textbf{Solution}

Suppose that the constant is $x$. Then $20+x, 50+x, 100+x$ is a \href{https://artofproblemsolving.com/wiki/index.php?title=Geometric\_progression}{geometric progression}, so $(20+x)(100+x) = (50+x)^2$. Expanding, we get $2000 + 120x + x^2 = 2500 + 100x + x^2$; therefore, $20x = 500$, so $x=25$.

 Now we can calculate our geometric progression to be $20+25, 50+25, 100+25 = 45, 75, 125$. Therefore, the common ratio is $\frac{125}{75} = \frac{75}{45} = \boxed{\frac{5}{3}}$, and our answer is $\boxed{\textbf{(A)}}$.

 


	\item (\textit{AHSME 1959}) The arithmetic mean (average) of a set of $50$ numbers is $38$. If two numbers, namely, $45$ and $55$, are discarded, the mean of the remaining set of numbers is:
$\textbf{(A)}\ 36.5 \qquad\textbf{(B)}\ 37\qquad\textbf{(C)}\ 37.2\qquad\textbf{(D)}\ 37.5\qquad\textbf{(E)}\ 37.52$


 
 \textbf{Answer: }D



\textbf{Solution}

Since the arithmetic mean of the $50$ numbers is $38$, their sum must be $50*38 = 1900$. After $45$ and $55$ are discarded, the sum decreases by $45 + 55 = 100$, so it must become $1900 - 100 = 1800$. But this means that the new mean of the remaining $50 - 2 = 48$ numbers must be $\frac{1800}{48} = 37.5$. Thusly, our answer is $\boxed{\textbf{(D)}}$, and we are done.

 


	\item (\textit{AHSME 1959}) Given the set $S$ whose elements are zero and the even integers, positive and negative. 
Of the five operations applied to any pair of elements: (1) addition (2) subtraction 
(3) multiplication (4) division (5) finding the arithmetic mean (average), those elements that only yield elements of $S$ are:
$\textbf{(A)}\ \text{all} \qquad\textbf{(B)}\ 1,2,3,4\qquad\textbf{(C)}\ 1,2,3,5\qquad\textbf{(D)}\ 1,2,3\qquad\textbf{(E)}\ 1,3,5$


 
 \textbf{Answer: }D



\textbf{Solution}

The first three listed operations, applied to even integers, all yield even integers trivially. As for the fourth operation (division) and fifth operation (average of pair), we can find pairs of even numbers that are mapped to odd integers, such as $\frac{2}{2} = 1$, and $\frac{0+2}{2} = 1$. Therefore, the operations that map even integers to even integers only are operations $1, 2, 3$, so our answer is $\boxed{\textbf{(D)}}$ and we are done.

 


	\item (\textit{AHSME 1959}) In a right triangle the square of the hypotenuse is equal to twice the product of the legs. One of the acute angles of the triangle is:
$\textbf{(A)}\ 15^{\circ} \qquad\textbf{(B)}\ 30^{\circ} \qquad\textbf{(C)}\ 45^{\circ} \qquad\textbf{(D)}\ 60^{\circ}\qquad\textbf{(E)}\ 75^{\circ}$


 
 \textbf{Answer: }C



\textbf{Solution 1}

WLOG, by scaling, that the hypotenuse has length 1. Let $\theta$ be an angle opposite from some leg. Then the two legs have length $\sin\theta$ and $\cos\theta$ respectively, so we have $2\sin\theta\cos\theta = 1^2$. From trigonometry, we know that this equation is true when $\theta = 45^{\circ}$, so our answer is $\boxed{\textbf{(C)}}$ and we are done.

 



\textbf{Solution 2 (Answer choices)}

Look at the options. Note that if $\textbf{(A)}$ is the correct answer, one of the acute angles of the triangle will measure to $15$ degrees. This implies that the other acute angle of the triangle would measure to be $75$ degrees, which would imply that $\textbf{(E)}$ is another correct answer. However, there is only one correct answer per question, so $\textbf{(A)}$ can't be a correct answer. Using a similar argument, neither $\textbf{(B)}$, $\textbf{(D)}$ , nor $\textbf{(E)}$ can be a correct answer. Since $\textbf{(C)}$ is the only answer choice left and there must be one correct answer, the answer must be $\boxed{\textbf{(C)}}$.

 


	\item (\textit{AHSME 1959}) The expression$\frac{x^2-3x+2}{x^2-5x+6}\div \frac{x^2-5x+4}{x^2-7x+12},$ when simplified is:
$\textbf{(A)}\ \frac{(x-1)(x-6)}{(x-3)(x-4)} \qquad\textbf{(B)}\ \frac{x+3}{x-3}\qquad\textbf{(C)}\ \frac{x+1}{x-1}\qquad\textbf{(D)}\ 1\qquad\textbf{(E)}\ 2$


 
 \textbf{Answer: }D



\textbf{Solution}

Factoring each of the binomials in the expression $\frac{x^2-3x+2}{x^2-5x+6}\div \frac{x^2-5x+4}{x^2-7x+12},$ will yield the result of 
 \[\frac{(x-2)(x-1)}{(x-3)(x-2)}\div \frac{(x-4)(x-1)}{(x-3)(x-4)},\]We can eliminate like terms to get $\frac {x-1}{x-3}\div \frac{x-1}{x-3}$, which, according to identity property, is equivalent to the answer $\boxed{\textbf{(D) }1}$.

 


	\item (\textit{AHSME 1959}) If $y=a+\frac{b}{x}$, where $a$ and $b$ are constants, and if $y=1$ when $x=-1$, and $y=5$ when $x=-5$, then $a+b$ equals:
$\textbf{(A)}\ -1 \qquad\textbf{(B)}\ 0\qquad\textbf{(C)}\ 1\qquad\textbf{(D)}\ 10\qquad\textbf{(E)}\ 11$


 
 \textbf{Answer: }E



\textbf{Solution}

Plugging in the x and y values, we obtain the following system of equations:
 \[\begin {cases} 1 = a - b \\ 5 = a - \frac{b}{5} \end {cases}\]We can then subtract the equations to obtain the equation $4 = 0.8b$, which works out to $b = 5$.

 Plugging that in to the original system of equations:
 \[\begin {cases} 1 = a - 5 \\ 5 = a - \frac{5}{5} \end {cases}\]It is easy to see that $a = 6$, and that $a+b$ is $\boxed{\textbf{(E) } 11}$.

 


	\item (\textit{AHSME 1959}) The arithmetic mean (average) of the first $n$ positive integers is:
$\textbf{(A)}\ \frac{n}{2} \qquad\textbf{(B)}\ \frac{n^2}{2}\qquad\textbf{(C)}\ n\qquad\textbf{(D)}\ \frac{n-1}{2}\qquad\textbf{(E)}\ \frac{n+1}{2}$


 
 \textbf{Answer: }E



\textbf{Solution}

The sum of the first $n$ positive integers is the same as the nth triangular number, which can be expressed as $\frac{(n)(n+1)}{2}$. Since the question is asking for the arithmetic mean, the whole sum is divided by $n$, thus giving us $\frac {\frac{(n)(n+1)}{2}}{n}$, which then simplifies to $\boxed{\textbf{(E) } \frac{n+1}{2}}$.

 


	\item (\textit{AHSME 1959}) With the use of three different weights, namely $1$ lb., $3$ lb., and $9$ lb., how many objects of different weights can be weighed, 
if the objects is to be weighed and the given weights may be placed in either pan of the scale?
$\textbf{(A)}\ 15 \qquad\textbf{(B)}\ 13\qquad\textbf{(C)}\ 11\qquad\textbf{(D)}\ 9\qquad\textbf{(E)}\ 7$


 
 \textbf{Answer: }B



\textbf{Solution}

The heaviest object that could be weighed with this set weighs $1 + 3 + 9 = 13$ lb., and we can weigh any positive integer weight at most that. This means that $13$ different objects could be weighed, so our answer is $\boxed{\textbf{(B)}}$ and we are done.

 


	\item (\textit{AHSME 1959}) It is given that $x$ varies directly as $y$ and inversely as the square of $z$, and that $x=10$ when $y=4$ and $z=14$. Then, when $y=16$ and $z=7$, $x$ equals:
$\textbf{(A)}\ 180\qquad \textbf{(B)}\ 160\qquad \textbf{(C)}\ 154\qquad \textbf{(D)}\ 140\qquad \textbf{(E)}\ 120$


 
 \textbf{Answer: }B



\textbf{Solution}

$x$ varies directly to $\frac{y}{z^2}$ (The inverse variation of y and the square of z)

 We can write the expression 

 $x = \frac{ky}{z^2}$

 Now we plug in the values of $x=10$ when $y=4$ and $z=14$. 

 This gives us $k=490$

 We can use this to find the value of $x$ when $y=4$ and $z=14$

 $x=\frac{490\cdot4}{14^2}$

 Simplifying this we get,

 $\fbox{\textbf{(B) } 160}$

 


	\item (\textit{AHSME 1959}) If $p$ is the perimeter of an equilateral $\triangle$ inscribed in a circle, the area of the circle is:
$\textbf{(A)}\ \frac{\pi p^2}{3} \qquad\textbf{(B)}\ \frac{\pi p^2}{9}\qquad\textbf{(C)}\ \frac{\pi p^2}{27}\qquad\textbf{(D)}\ \frac{\pi p^2}{81}\qquad\textbf{(E)}\ \frac{\pi p^2\sqrt3}{27}$


 
 \textbf{Answer: }C



\textbf{Solution}

A side length of the triangle is $\frac{p}3$. An altitude of the triangle, by 30-60-90 triangles, is $\frac{p\sqrt{3}}{6}$. Because all classical triangle centers coincide on an equilateral triangle, by centroid properties a circumradius is $\frac{p\sqrt{3}}{9}$. Finally, the area of the circumcircle is $\frac{\pi p^2}{27}\rightarrow\boxed{\textbf{C}}$.

 


	\item (\textit{AHSME 1959}) The line joining the midpoints of the diagonals of a trapezoid has length $3$. If the longer base is $97,$ then the shorter base is:
$\textbf{(A)}\ 94 \qquad\textbf{(B)}\ 92\qquad\textbf{(C)}\ 91\qquad\textbf{(D)}\ 90\qquad\textbf{(E)}\ 89$


 
 \textbf{Answer: }C



\textbf{Solution 1}


                    \begin{figure}[H]
                        centering
                        \includegraphics[width=0.5\textwidth]{Solution/AHSME_I_1959_22_Solution_Figure_1.png}
                    \end{figure}
                    

 Let $x$ be the length of the shorter base. Then:

 $3 = \frac{97-x}{2}$

 $6 = 97-x$ 

 $x = 91$

 Thus, our answer is $\boxed{\textbf{(C) }91}$.

 



\textbf{Solution 2}

Let the trapezoid be $ABCD$ with $\overline{AD} \parallel \overline{BC}$, with $M$ as the midpoint of $\overline{AC}$, and $N$ as the midpoint of $\overline{BD}$, as in the diagram. As in the first solution, let the shorter base of the trapezoid ($\overline{AD}$) have length $x$. Because $\overline{AD} \parallel \overline{BC}$, we can imagine shifting $\overline{AC}$ along $\overleftrightarrow{BC}$ by distance $x$ such that $A$ is at $D$, at which point we get the following triangle:

 
                    \begin{figure}[H]
                        centering
                        \includegraphics[width=0.5\textwidth]{Solution/AHSME_I_1959_22_Solution_Figure_2.png}
                    \end{figure}
                    

 Because $\overline{MN}$ is a \href{https://artofproblemsolving.com/wiki/index.php?title=Midpoint\#Midsegments}{midpoint connector} of $\triangle ABC$, $MN=\frac{1}{2}BC$, and so we have the equation $3+x=\frac{97+x}{2}$. Solving for $x$ yields $x=\boxed{\textbf{(C) }91}$.

 


	\item (\textit{AHSME 1959}) The set of solutions of the equation $\log_{10}\left( a^2-15a\right)=2$ consists of
$\textbf{(A)}\ \text{two integers } \qquad\textbf{(B)}\ \text{one integer and one fraction}\qquad \textbf{(C)}\ \text{two irrational numbers }\qquad\textbf{(D)}\ \text{two non-real numbers} \qquad\textbf{(E)}\ \text{no numbers, that is, the empty set}$


 
 \textbf{Answer: }A



\textbf{Solution}

Understand that $\log_{10}\left(x\right)=y$ can be expressed as $x = 10^y$.

 By applying this fact to $\log_{10}\left( a^2-15a\right)=2$, we get $a^2-15a = 10^2$

 We can use factoring methods to bring us to $(a-20)(a+5)=0$, which, as each binomial produces one integer solution, gives us a solution set of $\fbox{\textbf{(A) }two integers}$.

 


	\item (\textit{AHSME 1959}) A chemist has m ounces of salt water that is $m$\% salt. How many ounces of salt must he add to make a solution that is $2m$\% salt?


 $\textbf{(A)}\ \frac{m}{100+m} \qquad\textbf{(B)}\ \frac{2m}{100-2m}\qquad\textbf{(C)}\ \frac{m^2}{100-2m}\qquad\textbf{(D)}\ \frac{m^2}{100+2m}\qquad\textbf{(E)}\ \frac{2m}{100+2m}$


 
 \textbf{Answer: }C



\textbf{Solution}

There are $m$ oz of salt, and there is $m\times\frac{m}{100}=\frac{m^2}{100}$ oz of pure salt. We wish to add $k$ oz of pure salt so that $\frac{m^2/100+k}{m+k}=\frac{2m}{100}$. Solving this results in $k=\frac{m^2}{100-2m}\rightarrow\boxed{\textbf{C}}$

 


	\item (\textit{AHSME 1959}) The symbol $|a|$ means $+a$ if $a$ is greater than or equal to zero, and $-a$ if a is less than or equal to zero; the symbol $<$ means "less than"; 
the symbol $>$ means "greater than."
The set of values $x$ satisfying the inequality $|3-x|<4$ consists of all $x$ such that:
$\textbf{(A)}\ x^2<49 \qquad\textbf{(B)}\ x^2>1 \qquad\textbf{(C)}\ 1<x^2<49\qquad\textbf{(D)}\ -1<x<7\qquad\textbf{(E)}\ -7<x<1$


 
 \textbf{Answer: }D



\textbf{Solution}

The equation $|3-x| < 4$ can be solved by splitting it into two inequalities: $3-x<4$ and $3-x>-4$. The solutions to those inequalities are $x>-1$ and $x<7$, respectively. The common interval of those two inequalities is $\boxed{\textbf{(D)}\ -1<x<7}$.

 


	\item (\textit{AHSME 1959}) The base of an isosceles triangle is $\sqrt 2$. The medians to the legs intersect each other at right angles. The area of the triangle is:
$\textbf{(A)}\ 1.5 \qquad\textbf{(B)}\ 2\qquad\textbf{(C)}\ 2.5\qquad\textbf{(D)}\ 3.5\qquad\textbf{(E)}\ 4$


 
 \textbf{Answer: }A



\textbf{Solution}

Drop the altitude $AD$ from the vertex $A$ to base $BC$, and note that it intersects the centroid $M$. Then $BCM$ is an isosceles right triangle, so $MB=1$, so $MD=\frac{\sqrt{2}}{2}$. By centroid properties, $AD=3MD=\frac{3\sqrt2}{2}$. Then the area of $ABC$ is $\frac12\sqrt{2}\frac{3\sqrt2}{2}=1.5\rightarrow\boxed{\textbf{A}}$.

 


	\item (\textit{AHSME 1959}) Which one of the following is  not  true for the equation $ix^2-x+2i=0$, where $i=\sqrt{-1}$
$\textbf{(A)}\ \text{The sum of the roots is 2} \qquad \\ \textbf{(B)}\ \text{The discriminant is 9}\qquad \\ \textbf{(C)}\ \text{The roots are imaginary}\qquad \\ \textbf{(D)}\ \text{The roots can be found using the quadratic formula}\qquad \\ \textbf{(E)}\ \text{The roots can be found by factoring, using imaginary numbers}$


 
 \textbf{Answer: }A



\textbf{Solution}

By \href{https://artofproblemsolving.com/wiki/index.php?title=Vieta\%27s\_Formulas}{Vieta's Formulas}, the sum of the roots is $\frac{1}{i} = -i$, which is not $2$, so the solution is $\fbox{\textbf{(A)}}$.

 


	\item (\textit{AHSME 1959}) In triangle $ABC$, $AL$ bisects angle $A$, and $CM$ bisects angle $C$. Points $L$ and $M$ are on $BC$ and $AB$, respectively. The sides of $\triangle ABC$ are $a$, $b$, and $c$. Then $\frac{AM}{MB}=k\frac{CL}{LB}$ where $k$ is:
$\textbf{(A)}\ 1 \qquad\textbf{(B)}\ \frac{bc}{a^2}\qquad\textbf{(C)}\ \frac{a^2}{bc}\qquad\textbf{(D)}\ \frac{c}{b}\qquad\textbf{(E)}\ \frac{c}{a}$


 
 \textbf{Answer: }E



\textbf{Solution}


                    \begin{figure}[H]
                        centering
                        \includegraphics[width=0.5\textwidth]{Solution/AHSME_I_1959_28_Solution_Figure_1.png}
                    \end{figure}
                    

 By the \href{https://artofproblemsolving.com/wiki/index.php?title=Angle\_Bisector\_Theorem}{Angle Bisector Theorem}, $\frac{AM}{AB}=\frac{AC}{BC}$ and $\frac{CL}{LB}=\frac{AC}{AB}$, so by rearranging the given equation and noting $AB=c$ and $BC=a$, $k=\frac{c}{a}\rightarrow\boxed{\textbf{(E)}}$.

 


	\item (\textit{AHSME 1959}) On a examination of $n$ questions a student answers correctly $15$ of the first $20$. Of the remaining questions he answers one third correctly. 
All the questions have the same credit. If the student's mark is 50\%, how many different values of $n$ can there be?
$\textbf{(A)}\ 4 \qquad\textbf{(B)}\ 3\qquad\textbf{(C)}\ 2\qquad\textbf{(D)}\ 1\qquad\textbf{(E)}\ \text{the problem cannot be solved}$


 
 \textbf{Answer: }D



\textbf{Solution}

To calculate the student's score in terms of $n$, you can write the following equation:

 $\frac{\frac{n-20}{3} + 15}{n} = \frac{1}{2}$. Simplify to get $n=50$, so there is $\boxed{\textbf{(D) }1}$ solution.

 


	\item (\textit{AHSME 1959}) $A$ can run around a circular track in $40$ seconds. $B$, running in the opposite direction, meets $A$ every $15$ seconds. 
What is $B$'s time to run around the track, expressed in seconds?
$\textbf{(A)}\ 12\frac12 \qquad\textbf{(B)}\ 24\qquad\textbf{(C)}\ 25\qquad\textbf{(D)}\ 27\frac12\qquad\textbf{(E)}\ 55$


 
 \textbf{Answer: }B



\textbf{Solution}

In $15$ seconds, A will complete $\frac{3}{8}$ of the track. This means that B will complete $\frac{5}{8}$ of the track in $15$ seconds, meaning that to complete the whole track (making the fraction 1), it will take $\frac{8}{5} \cdot 15 = 24$ seconds. So the answer is $\boxed{\textbf{(B) }24}$.

 


	\item (\textit{AHSME 1959}) A square, with an area of $40$, is inscribed in a semicircle. The area of a square that could be inscribed in the entire circle with the same radius, is:
$\textbf{(A)}\ 80 \qquad\textbf{(B)}\ 100\qquad\textbf{(C)}\ 120\qquad\textbf{(D)}\ 160\qquad\textbf{(E)}\ 200$


 
 \textbf{Answer: }B



\textbf{Solution}


                    \begin{figure}[H]
                        centering
                        \includegraphics[width=0.5\textwidth]{Solution/AHSME_I_1959_31_Solution_Figure_1.png}
                    \end{figure}
                    

 Let the points be labeled as in the diagram, with $O$ being the center of the circle. Because we know that the small square has an area of $40$, it must have a side length of $\sqrt{40}=2\sqrt{10}$. It is simple to prove that $O$ is the midpoint of the bottom side of the small square, so $OB=\frac{2\sqrt{10}}{2}=\sqrt{10}$. By the \href{https://artofproblemsolving.com/wiki/index.php?title=Pythagorean\_Theorem}{Pythagorean Theorem}, $AO=\sqrt{50}=5\sqrt{2}$, which is the radius of the circle. Thus, $OC=5\sqrt{2}$, and so the diagonal $\overline{CD}$ of the big square has length $10\sqrt{2}$. Thus, the big square has side length $\frac{10\sqrt{2}}{\sqrt{2}}=10$, and, consequently, it has an area of $\boxed{\textbf{(B) }100}$.

 


	\item (\textit{AHSME 1959}) The length $l$ of a tangent, drawn from a point $A$ to a circle, is $\frac43$ of the radius $r$. The (shortest) distance from A to the circle is:
$\textbf{(A)}\ \frac{1}{2}r \qquad\textbf{(B)}\ r\qquad\textbf{(C)}\ \frac{1}{2}l\qquad\textbf{(D)}\ \frac23l \qquad\textbf{(E)}\ \text{a value between r and l.}$


 
 \textbf{Answer: }C



\textbf{Solution}


                    \begin{figure}[H]
                        centering
                        \includegraphics[width=0.5\textwidth]{Solution/AHSME_I_1959_32_Solution_Figure_1.png}
                    \end{figure}
                    

 Let the circle have center $O$, let the point of tangency be point $T$, and let $B$ be the intersection of $\overline{OA}$ with the circle, as in the diagram. By the definitions of a circle and a tangent to a circle, we know that $OB=OT=r$ and $\overline{OT} \perp \overline{TA}$. By the \href{https://artofproblemsolving.com/wiki/index.php?title=Pythagorean\_Theorem}{Pythagorean Theorem}, $OA=\sqrt{\frac{25r^2}{9}}=\frac{5}{3}r$. Because the shortest segment from an external point to a circle lies on the line connecting that point to the center of the circle, our desired distance is $AB$. Because $OA=\frac{5}{3}r$ and $OB=r$, $AB=\frac{5}{3}r-r=\frac{2}{3}r=\frac{(4/3)r}{2}=\boxed{\textbf{(C) }\frac{1}{2}l}$.

 


	\item (\textit{AHSME 1959}) A harmonic progression is a sequence of numbers such that their reciprocals are in arithmetic progression.
Let $S_n$ represent the sum of the first $n$ terms of the harmonic progression; for example $S_3$ represents the sum of 
the first three terms. If the first three terms of a harmonic progression are $3,4,6$, then:
$\textbf{(A)}\ S_4=20 \qquad\textbf{(B)}\ S_4=25\qquad\textbf{(C)}\ S_5=49\qquad\textbf{(D)}\ S_6=49\qquad\textbf{(E)}\ S_2=\frac{1}2 S_4$


 
 \textbf{Answer: }B



\textbf{Solution}

From the problem, we know that $\frac {1} {3}$,$\frac {1} {4}$,$\frac {1} {6}$, ... are in arithmetic progression.

 Then, the common difference $d$ of this arithmetic progression is $\frac {1} {4}$ $-$ $\frac {1} {3}$ $=$ $\frac {1} {6}$ $-$ $\frac {1} {4}$ $=$ $-$ $\frac {1} {12}$

 Thus, the fourth term of the arithmetic progression is $\frac{1}{12}$.

 So, the fourth term of the harmonic progression is $12$. 

 Now, we can see that $S_4$ $=$ $3$ $+$ $4$ $+$ $6$ $+$ $12$ $=$ $25$, which corresponds to answer $\fbox{\textbf{(B)}}$.

 


	\item (\textit{AHSME 1959}) Let the roots of $x^2-3x+1=0$ be $r$ and $s$. Then the expression $r^2 + s^2$ is:
$\textbf{(A)}\ \text{a positive integer} \qquad\textbf{(B)}\ \text{a positive fraction greater than 1}\qquad\textbf{(C)}\ \text{a positive fraction less than 1}\qquad\textbf{(D)}\ \text{an irrational number}\qquad\textbf{(E)}\ \text{an imaginary number}$


 
 \textbf{Answer: }A



\textbf{Solution}

You may recognize that $r^2+s^2$ can be written as $(r+s)^2-2rs$. Then, by \href{https://artofproblemsolving.com/wiki/index.php?title=Vieta\%27s\_formulas}{Vieta's formulas}, 
 \[r+s=-(-3)=3\]and 
 \[rs=1.\]Therefore, plugging in the values for $r+s$ and $rs$, 
 \[(r+s)^2-2rs=(3)^2-2(1)=9-2=7.\]Hence, we can say that the expression $r^2+s^2$ is $\boxed{\textbf{(A)}\ \text{a positive integer.}}$

 


	\item (\textit{AHSME 1959}) The symbol $\ge$ means "greater than or equal to"; the symbol $\le$ means "less than or equal to".
In the equation $(x-m)^2-(x-n)^2=(m-n)^2; m$ is a fixed positive number, and $n$ is a fixed negative number. The set of values x satisfying the equation is:
$\textbf{(A)}\ x\ge 0 \qquad\textbf{(B)}\ x\le n\qquad\textbf{(C)}\ x=0\qquad\textbf{(D)}\ \text{the set of all real numbers}\qquad\textbf{(E)}\ \text{none of these}$


 
 \textbf{Answer: }E



\textbf{Solution}

Applying the \href{https://artofproblemsolving.com/wiki/index.php?title=Difference\_of\_squares}{difference of squares} technique on this problem, we can see that 
 \[(x-m)^2-(x-n)^2=((x-m)+(x-n)) \cdot ((x-m)-(x-n)),\] so 
 \[(x-m)^2-(x-n)^2=((x-m)+(x-n)) \cdot ((x-m)-(x-n))=(m-n)^2.\] Simplifying gives us
 \[(2x-m-n)\cdot(n-m)=(m-n)^2.\]Negating $n-m$ creates:
 \[-(2x-m-n)\cdot(m-n)=(m-n)^2.\]Dividing by $m-n$, 
 \[-(2x-m-n)=m-n\] 
 \[-2x+m+n=m-n\] 
 \[-2x+n=-n\] 
 \[-2x=-2n\] 
 \[x=n\]Lastly, since $n$ is a fixed negative number, $x$ must also be a fixed negative number, so $x<0$. Since this answer is not $A, B, C,$ or $D$, the solution must be $\fbox{\textbf{(E) }none of these}$.

 


	\item (\textit{AHSME 1959}) The base of a triangle is $80$, and one side of the base angle is $60^\circ$. The sum of the lengths of the other two sides is $90$. The shortest side is:
$\textbf{(A)}\ 45 \qquad\textbf{(B)}\ 40\qquad\textbf{(C)}\ 36\qquad\textbf{(D)}\ 17\qquad\textbf{(E)}\ 12$


 
 \textbf{Answer: }D



\textbf{Solution}


                    \begin{figure}[H]
                        centering
                        \includegraphics[width=0.5\textwidth]{Solution/AHSME_I_1959_36_Solution_Figure_1.png}
                    \end{figure}
                    

 Let the triangle be $\triangle ABC$ with $\measuredangle BAC=60^{\circ}$ and $AC=80$, as in the diagram. Let $AB=x$. from the problem, we know that $BC=90-x$. Now, we can apply the \href{https://artofproblemsolving.com/wiki/index.php?title=Law\_of\_Cosines}{Law of Cosines} on $\triangle ABC$ to solve for $x$:\begin{align*}(90-x)^2 &= 6400+x^2-160\cos(60^{\circ}) \\x^2-180x+8100 &= 6400+x^2-80x \\-100x &= -1700 \\x &= 17\end{align*}
Because the problem asks for the shortest side of the triangle and $90-x=73>17$, our answer is $\boxed{\textbf{(D) }17}$.

 


	\item (\textit{AHSME 1959}) When simplified the product $\left(1-\frac13\right)\left(1-\frac14\right)\left(1-\frac15\right)\cdots\left(1-\frac1n\right)$ becomes:
$\textbf{(A)}\ \frac1n \qquad\textbf{(B)}\ \frac2n\qquad\textbf{(C)}\ \frac{2(n-1)}{n}\qquad\textbf{(D)}\ \frac{2}{n(n+1)}\qquad\textbf{(E)}\ \frac{3}{n(n+1)}$


 
 \textbf{Answer: }B



\textbf{Solution}

The product  $\left(1-\frac13\right)\left(1-\frac14\right)\left(1-\frac15\right)\cdots\left(1-\frac1n\right)$ can be simplified to $\frac{2*3*4\cdots*(n-2)*(n-1)}{3*4*5\cdots*(n-1)*n}$, which, through \href{https://artofproblemsolving.com/wiki/index.php?title=Telescoping\_series}{telescoping}, ultimately works out to $\boxed{\textbf{(B) }\frac{2}{n}}$.

 


	\item (\textit{AHSME 1959}) If $4x+\sqrt{2x}=1$, then $x$:
$\textbf{(A)}\ \text{is an integer} \qquad\textbf{(B)}\ \text{is fractional}\qquad\textbf{(C)}\ \text{is irrational}\qquad\textbf{(D)}\ \text{is imaginary}\qquad\textbf{(E)}\ \text{may have two different values}$


 
 \textbf{Answer: }B



\textbf{Solution}

Subtract $4x$ from both sides of the initial equation and solve for $x$:\begin{align*}\sqrt{2x} &= 1-4x \\2x &= 1-8x+16x^2 \\16x^2-10x+1 &= 0 \\(8x-1)(2x-1) &= 0\end{align*}


 It may look like we have two different solutions for $x$, but plugging in $x=\frac{1}{2}$ into the original equation reveals that it is \href{https://artofproblemsolving.com/wiki/index.php?title=Extraneous\_solutions}{extraneous}. Thus, $x=\frac{1}{8}$, which is only consistent with answer choice $\fbox{\textbf{(B)}}$.

 


	\item (\textit{AHSME 1959}) Let S be the sum of the first nine terms of the sequence $x+a, x^2+2a, x^3+3a, \cdots.$
Then S equals:
$\textbf{(A)}\ \frac{50a+x+x^8}{x+1} \qquad\textbf{(B)}\ 50a-\frac{x+x^{10}}{x-1}\qquad\textbf{(C)}\ \frac{x^9-1}{x+1}+45a\qquad\textbf{(D)}\ \frac{x^{10}-x}{x-1}+45a\qquad\textbf{(E)}\ \frac{x^{11}-x}{x-1}+45a$


 
 \textbf{Answer: }D



\textbf{Solution}

We know $x+x^2\dots+x^9=x(1+x\dots+x^8)=x\frac{x^9-1}{x-1}=\frac{x^{10}-x}{x-1}$. The only answer with this term is $\boxed{\textbf{D}}$.

 


	\item (\textit{AHSME 1959}) In $\triangle ABC$, $BD$ is a median. $CF$ intersects $BD$ at $E$ so that $\overline{BE}=\overline{ED}$. Point $F$ is on $AB$. Then, if $\overline{BF}=5$, 
$\overline{BA}$ equals:
$\textbf{(A)}\ 10 \qquad\textbf{(B)}\ 12\qquad\textbf{(C)}\ 15\qquad\textbf{(D)}\ 20\qquad\textbf{(E)}\ \text{none of these}$


 
 \textbf{Answer: }C



\textbf{Solution 1}


                    \begin{figure}[H]
                        centering
                        \includegraphics[width=0.5\textwidth]{Solution/AHSME_I_1959_40_Solution_Figure_1.png}
                    \end{figure}
                    

 Draw $\overline{DG} \parallel \overline{FC}$ with $G$ on $\overline{AB}$. We know that $GF = BF = 5$, since $\triangle BFE \sim \triangle BGD$.

 Likewise, since $\triangle ADG \sim \triangle ACF$, we know that $AG=5$.

 Thus, $AF=AG+GF+FB=5+5+5=15$, which is answer $\fbox{\textbf{(C)}}$.

 


 



\textbf{Solution 2}

Let $AD=DC=x$ and $BE=ED=y$. By \href{https://artofproblemsolving.com/wiki/index.php?title=Menelaus\%27\_Theorem}{Menelaus' Theorem} on $\triangle BEF$ and $\overleftrightarrow{AC}$, we know the following:\begin{align*}\frac{BD}{ED}*\frac{EC}{FC}*\frac{FA}{BA} &= 1 \\\frac{2y}{y}*\frac{EC}{FC}*\frac{FA}{5+FA} &= 1 \\\frac{EC}{FC}*\frac{FA}{5+FA} &= \frac{1}{2}\end{align*}
By applying Menelaus again on $\triangle CDE$ and $\overleftrightarrow{AB}$, we see that:\begin{align*}\frac{CA}{DA}*\frac{DB}{EB}*\frac{EF}{CF} &= 1 \\\frac{2x}{x}*\frac{2y}{y}*\frac{CF-EC}{FC} &= 1 \\1-\frac{EC}{FC} &= \frac{1}{4} \\\frac{EC}{FC} &= \frac{3}{4}\end{align*}
Substituting $\frac{3}{4}$ for $\frac{EC}{FC}$ into the previous equation, we can now solve for $FA$:\begin{align*}\frac{EC}{FC}*\frac{FA}{5+FA} &= \frac{1}{2} \\\frac{3}{4}*\frac{FA}{5+FA} &= \frac{1}{2} \\\frac{FA}{5+FA} &= \frac{2}{3} \\3FA &= 10+2FA \\FA &= 10\end{align*}
Because $AB=AF+FB$, $AB=5+10=\boxed{\textbf{(C) } 15}$.

 


	\item (\textit{AHSME 1959}) On the same side of a straight line three circles are drawn as follows: a circle with a radius of $4$ inches is tangent to the line, the other 
two circles are equal, and each is tangent to the line and to the other two circles. The radius of the equal circles is:
$\textbf{(A)}\ 24 \qquad\textbf{(B)}\ 20\qquad\textbf{(C)}\ 18\qquad\textbf{(D)}\ 16\qquad\textbf{(E)}\ 12$


 
 \textbf{Answer: }D



\textbf{Solution}


                    \begin{figure}[H]
                        centering
                        \includegraphics[width=0.5\textwidth]{Solution/AHSME_I_1959_41_Solution_Figure_1.png}
                    \end{figure}
                    

 Let the radius of the two large circles be $r$. Also, let the center of the small circle be $A$ and the centers of the two large circles be $B$ and $C$, as in the diagram. Further, let the large circles intersect at $F$, let the circle centered at $B$ be tangent to the line at $D$, and let the circle centered at $C$ be tangent to the line at $E$. Because $BD=CE=r$ and $\overline{BD},\overline{CE} \perp \overleftrightarrow{DE}$, $\overline{BD} \parallel \overleftrightarrow{DE}$. Because the two large circles are tangent at $F$, $F$ is on the segment connecting their centers, $\overline{BC}$. Thus, $F$ is a distance $r$ from $\overleftrightarrow{DE}$. Because the small circle (with radius $4$) is tangent to the line, $AF=r-4$. Also, because the small circle and the circle centered at $B$ are tangent, $AB=r+4$. Because $BF=r$, by the \href{https://artofproblemsolving.com/wiki/index.php?title=Pythagorean\_Theorem}{Pythagorean Theorem}, we have the following equation:\begin{align*}BF^2+AF^2 &= AB^2 \\r^2+(r-4)^2 &= (r+4)^2 \\r^2+r^2-8r+16 &= r^2+8r+16 \\r^2-16r &= 0 \\r(r-16) &= 0\end{align*}
Because $r \neq 0$, we are left with $r=16$. Thus, the radius of the two large circles is $\boxed{\textbf{(D) }16}$.

 


	\item (\textit{AHSME 1959}) Given three positive integers $a,b,$ and $c$. Their greatest common divisor is $D$; their least common multiple is $M$. 
Then, which two of the following statements are true?
$\text{(1)}\ \text{the product MD cannot be less than abc} \qquad \\ \text{(2)}\ \text{the product MD cannot be greater than abc}\qquad \\ \text{(3)}\ \text{MD equals abc if and only if a,b,c are each prime}\qquad \\ \text{(4)}\ \text{MD equals abc if and only if a,b,c are each relatively prime in pairs} \text{ (This means: no two have a common factor greater than 1.)}$
$\textbf{(A)}\ 1,2 \qquad\textbf{(B)}\ 1,3\qquad\textbf{(C)}\ 1,4\qquad\textbf{(D)}\ 2,3\qquad\textbf{(E)}\ 2,4$


 
 \textbf{Answer: }E



\textbf{Solution}

Because $1\times2\times4>1\times4$, 1 is false. Because $1\times1\times1=1\times1$, 3 is false. It follows that the answer is $\boxed{\textbf{E}}$.

 


	\item (\textit{AHSME 1959}) The sides of a triangle are $25,39$, and $40$. The diameter of the circumscribed circle is:
$\textbf{(A)}\ \frac{133}{3}\qquad\textbf{(B)}\ \frac{125}{3}\qquad\textbf{(C)}\ 42\qquad\textbf{(D)}\ 41\qquad\textbf{(E)}\ 40$


 
 \textbf{Answer: }B



\textbf{Solution}

The semiperimeter $s$ of the triangle is $\frac{25+39+40}{2}=52$. Therefore, by \href{https://artofproblemsolving.com/wiki/index.php?title=Heron\%27s\_Formula}{Heron's Formula}, we can find the area of the triangle as follows:\begin{align*}A &= \sqrt{52(52-25)(52-39)(52-40)} \\&= \sqrt{52(27)(13)(12)} \\&= \sqrt{13*4(9*3)(13)(4*3)} \\&= 13*4*3*3 \\&= 468\end{align*}
Also, we know that the area of the triangle is $\frac{abc}{4R}$, where $a$, $b$, and $c$ are the sides of the triangle and $R$ is the triangle's \href{https://artofproblemsolving.com/wiki/index.php?title=Circumradius}{circumradius}. Thus, we can equate this expression for the area with $468$ to solve for $R$:\begin{align*}\frac{(25)(39)(40)}{4R} &= 468 = 39*12 \\\frac{25*40}{12} &= 4R \\R &= \frac{250}{12} = \frac{125}{6}\end{align*}
Because the problem asks for the diameter of the circumcircle, our answer is $2R=\boxed{\textbf{(B) }\frac{125}{3}}$.

 


	\item (\textit{AHSME 1959}) The roots of $x^2+bx+c=0$ are both real and greater than $1$. Let $s=b+c+1$. Then $s$:
$\textbf{(A)}\ \text{may be less than zero}\qquad\textbf{(B)}\ \text{may be equal to zero}\qquad$


 $\textbf{(C)} \text{ must be greater than zero}\qquad\textbf{(D)}\ \text{must be less than zero}\qquad  \textbf{(E)}\text{ must be between -1 and +1}$


 
 \textbf{Answer: }C



\textbf{Solution}

Let the roots of the quadratic be $p$ and $q$. Then, by \href{https://artofproblemsolving.com/wiki/index.php?title=Vieta\%27s\_Formulas}{Vieta's Formulas}, $b=-(p+q)$ and $c=pq$. By substituting these values of $b$ and $c$ into our expression for $s$, we see that $s=pq-p-q+1$. By \href{https://artofproblemsolving.com/wiki/index.php?title=SFFT}{SFFT}, $s=(p-1)(q-1)$. From the problem, we know that $p$ and $q$ are both greater than $1$, so $(p-1)$ and $(q-1)$ are necessarily positive. Thus, $s$, the product of these two positive terms, $\fbox{\textbf{(C) }must be greater than zero}$.

 


	\item (\textit{AHSME 1959}) If $\left(\log_3 x\right)\left(\log_x 2x\right)\left( \log_{2x} y\right)=\log_{x}x^2$, then $y$ equals:


 $\textbf{(A)}\ \frac92\qquad\textbf{(B)}\ 9\qquad\textbf{(C)}\ 18\qquad\textbf{(D)}\ 27\qquad\textbf{(E)}\ 81$


 
 \textbf{Answer: }B



\textbf{Solution}

From the properties of \href{https://artofproblemsolving.com/wiki/index.php?title=Logarithms}{logarithms}, we can simplify the equation and solve for $y$:\begin{align*}(\log_3 x)(\log_x 2x)(\log_{2x} y) &= \log_{x}x^2 \\(\log_3 2x)(\log_{2x} y) &= 2\log_x x \\\log_3 y &= 2 \\y &= 3^2 \\y &= 9\end{align*}
Thus, our answer is $\boxed{\textbf{(B) }9}$.

 


	\item (\textit{AHSME 1959}) A student on vacation for $d$ days observed that (1) it rained $7$ times, morning or afternoon (2) when it rained in the afternoon, 
it was clear in the morning (3) there were five clear afternoons (4) there were six clear mornings. Then $d$ equals:
$\textbf{(A)}\ 7\qquad\textbf{(B)}\ 9\qquad\textbf{(C)}\ 10\qquad\textbf{(D)}\ 11\qquad\textbf{(E)}\ 12$


 
 \textbf{Answer: }B



\textbf{Solution}

From the given information, we know that there were $7+5+6=18$ halves of the day (either afternoons or mornings), so there were $\frac{18}{2}=9$ days in total, which is answer choice $\boxed{\textbf{(B)}}$.

 


	\item (\textit{AHSME 1959}) Assume that the following three statements are true:
(I). All freshmen are human. (II). All students are human. (III). Some students think.


 Given the following four statements:


 $\textbf{(1)}\ \text{All freshmen are students.}\qquad \\ \textbf{(2)}\ \text{Some humans think.}\qquad \\ \textbf{(3)}\ \text{No freshmen think.}\qquad \\ \textbf{(4)}\ \text{Some humans who think are not students.}$


 Those which are logical consequences of I, II, and III are:


 $\textbf{(A)}\ 2\qquad\textbf{(B)}\ 4\qquad\textbf{(C)}\ 2,3\qquad\textbf{(D)}\ 2,4\qquad\textbf{(E)}\ 1,2$


 
 \textbf{Answer: }A



\textbf{Solution}

All that we know about the freshmen is that they are all human (by statement I), so we can not prove assertion (1).

 Because all students are human (by statement II) and some students think (by statement III), some humans think, which proves assertion (2).

 Again, concerning the freshmen, all we know is that they are human. Thus, we cannot prove assertion (3). 

 Concerning those who are not students, all that we can say is that some freshmen or humans might not be students. This information is not sufficient to prove assertion (4).

 Thus, because we can only prove assertion (2), our answer is $\fbox{\textbf{(A)}}$.

 


	\item (\textit{AHSME 1959}) Given the polynomial $a_0x^n+a_1x^{n-1}+\cdots+a_{n-1}x+a_n$, where $n$ is a positive integer or zero, and $a_0$ is a positive integer. 
The remaining $a$'s are integers or zero. Set $h=n+a_0+|a_1|+|a_2|+\cdots+|a_n|$. [See example 25 for the meaning of $|x|$.] 
The number of polynomials with $h=3$ is:
$\textbf{(A)}\ 3\qquad\textbf{(B)}\ 5\qquad\textbf{(C)}\ 6\qquad\textbf{(D)}\ 7\qquad\textbf{(E)}\ 9$


 
 \textbf{Answer: }B



\textbf{Solution}

We perform \href{https://artofproblemsolving.com/wiki/index.php?title=Casework}{casework} by the value of $n$, the \href{https://artofproblemsolving.com/wiki/index.php?title=Degree}{degree} of our \href{https://artofproblemsolving.com/wiki/index.php?title=Polynomial}{polynomial} $a_0x^n+a_1x^{n-1}+\cdots+a_{n-1}x+a_n$.

 Case $n = 0$: In this case we are forced to set $a_0 = 3$. This contributes $1$ possibility.

 Case $n = 1$: In this case we must have $a_0 + |a_1| = 2$, so our polynomial could be $1 + 1x, 0 + 2x, -1 + 1x$. This contributes $3$ possibilities.

 Case $n = 2$: In this case we must have $a_0 + |a_1| + |a_2| = 1$. However, because $a_0$ must be positive, it has to be $1$, so our polynomial can only be $0 + 0x + 1x^2$. This contributes $1$ possibility.

 Case $n\geq 3$: This case is impossible because $h = n+a_0+|a_1|+|a_2|+\cdots+|a_n|\geq n + a_0\geq 3 + 1 = 4 > 3$, so it contributes $0$ possibilities.

 Adding the results from all four cases, we find that there are $1 + 3 + 1 + 0 = 5$ possibilities in total, so our answer is $\boxed{(B)}$.

 


	\item (\textit{AHSME 1959}) For the infinite series $1-\frac12-\frac14+\frac18-\frac{1}{16}-\frac{1}{32}+\frac{1}{64}-\frac{1}{128}-\cdots$ let $S$ be the (limiting) sum. Then $S$ equals:
$\textbf{(A)}\ 0\qquad\textbf{(B)}\ \frac27\qquad\textbf{(C)}\ \frac67\qquad\textbf{(D)}\ \frac{9}{32}\qquad\textbf{(E)}\ \frac{27}{32}$


 
 \textbf{Answer: }B



\textbf{Solution}

Notice that we can arrange the sequence like so: 
 \begin{align*}1-\frac12-\frac14+\frac18-\frac{1}{16}-\frac{1}{32}+\frac{1}{64}-\frac{1}{128}-\cdots  &= 1-\frac{5}{8}-\frac{5}{64}-\frac{5}{512}-\dots \\ &= 1-\left(\sum_{i=1}^{\infty}\frac{5}{8^n}\right) \\ &= 1-\frac{5}{7} \text{ (by the convergence of a geometric series)} \\ &=\frac{2}{7}\end{align*}
hence our answer is $\fbox{B}$

 


	\item (\textit{AHSME 1959}) A club with $x$ members is organized into four committees in accordance with these two rules:


 
 \[\text{(1)}\ \textup{Each member belongs to two and only two committees}\qquad \\ \text{(2)}\ \textup{Each pair of committees has one and only one member in common}\]


 Then $x$:


 $\textbf{(A)} \ \textup{cannot be determined} \qquad \\ \textbf{(B)} \ \textup{has a single value between 8 and 16} \qquad \\ \textbf{(C)} \ \textup{has two values between 8 and 16} \qquad \\ \textbf{(D)} \ \textup{has a single value between 4 and 8} \qquad \\ \textbf{(E)} \ \textup{has two values between 4 and 8}$


 
 \textbf{Answer: }D



\textbf{Solution}

We can label each of the $x$ members with the pair of committees that they belong to, which is clearly valid due to rule (1). Then, by rule (2), for each pair of committees, there is exactly one member labeled with that pair. But since we have four committees, there must be $x = \binom{4}{2} = 6$ members in total. Thusly our choice is $\boxed{\textbf{(D)}}$, and we are done.

 


\end{enumerate}